\documentclass[14pt,a4paper]{extarticle}
\usepackage[T2A]{fontenc}
\usepackage[utf8]{inputenc}
\usepackage[english,russian]{babel}
\usepackage{lmodern}
\usepackage{geometry}
\usepackage{array}
\usepackage{tabularx}
\usepackage{longtable}
\usepackage{ltablex}
\usepackage{booktabs}
\usepackage{hyperref}
\usepackage{indentfirst}
\usepackage{microtype}
\usepackage{graphicx}
\usepackage{setspace}
\usepackage{listings}
\usepackage{graphicx}
\usepackage{caption}
\usepackage{float}
\usepackage{tocloft}
\usepackage{titlesec}
\usepackage{newtxtext,newtxmath}

\titleformat{\section}        % форматируем \section
  {\normalfont\bfseries}      % стиль: обычный + жирный
  {\thesection.}              % номер секции с точкой
  {2mm}                       % расстояние между номером и текстом
  {}                           % перед заголовком можно добавить что-то
\titlespacing*{\section}{1.5cm}{*3}{*1}  % {отступ слева}{отступ сверху}{отступ снизу}

\captionsetup{justification=centering}
\renewcommand{\ttdefault}{cmtt}
\keepXColumns
\setlength{\parindent}{3em}
\newgeometry{left=2.5cm,right=2.5cm,top=3cm,bottom=2.5cm}

\renewcommand{\contentsname}{\hspace*{4em}{\smallСодержание}}
\setlength{\cftbeforesecskip}{0.5em}
\renewcommand{\cftsecleader}{\cftdotfill{\cftdotsep}}
\renewcommand{\cftsecaftersnum}{.\hspace{2mm}}

\title{Отчёт по "Программно-аппаратные средства интеллектуальных систем"}
\author{Выполнил студент группы М-РИТ-25-1, Станиславчук Сергей}

\begin{document}

\setstretch{1.5}
\noindent
\begin{minipage}[t]{0.01\textwidth}
    \includegraphics{~/programmer/++Programmer/5_1/LSTU.jpg}
\end{minipage}%
\begin{minipage}[c]{1.1\linewidth}
    \centering
    {\small
    МИНИСТЕРСТВО НАУКИ И ВЫСШЕГО ОБРАЗОВАНИЯ РФ \\[1em]
    ФЕДЕРАЛЬНОЕ ГОСУДАРСТВЕННОЕ БЮДЖЕТНОЕ \\
    ОБРАЗОВАТЕЛЬНОЕ УЧРЕЖДЕНИЕ ВЫСШЕГО ОБРАЗОВАНИЯ \\
    «ЛИПЕЦКИЙ ГОСУДАРСТВЕННЫЙ ТЕХНИЧЕСКИЙ УНИВЕРСИТЕТ»
    }
\end{minipage}%

\vspace{3em}

\noindent
\hspace{1em}%
\begin{minipage}[t]{0.8\textwidth}
\textbf{Институт} \hspace{2em} компьютерных наук \\
\textbf{Кафедра} \hspace{2em} автоматизированных систем управления
\end{minipage}


\vspace{6em}

\begin{center}
Сочинение по дисциплине \\ <<Планирование карьеры>>
\end{center}

\vspace{6em}

\noindent
\begin{tabular*}{\textwidth}{@{\extracolsep{\fill}} l c r }
    Студент М-РИТ-25-1 & \makebox[4cm]{\hrulefill} & Станиславчук С.М. \\[2em]
Руководитель & \\
Доцент & \makebox[4cm]{\hrulefill} & Триханова И.Е. \\
\end{tabular*}

\vfill
\begin{center}
    Липецк, 2026
\end{center}

\thispagestyle{empty}

\newpage
\tableofcontents
\newpage

\section{Эффект бабочки в строке кода}

Есть особая, почти сакральная тишина, которая наступает глубокой ночью, когда
последняя ошибка в консоли исчезает, и сцена, над которой ты бился несколько
дней, вдруг оживает. Именно в такие моменты понимаешь: разработка игр — это не
профессия и даже не искусство. Это способ мышления. Я нахожусь в самом начале
этого пути, и если смотреть с высоты птичьего полета, моя траектория может
показаться прагматичной: первый курс магистратуры, учеба, планы на будущее. Но
внутри этой рациональной оболочки уже зреет хаос творческих амбиций, который
трудно измерить стандартными карьерными метриками.

Сейчас мой главный соратник и оппонент — это личный проект на Unity. У меня нет
команды, нет дедлайнов от издателя и нет бюджета. У меня есть только чистый
лист, который я заполняю архитектурой игровых систем. Это, пожалуй, самый
важный этап в карьере разработчика, который часто недооценивают. Здесь, в
вакууме личной лаборатории, я имею роскошь ошибаться. Я могу переписать систему
инвентаря трижды, убить персонажа, который не вписывается в логику мира, и
позволить себе эксперименты, не грозящие срывом релиза. Это та самая «игра в
песочнице», которая формирует инженерную интуицию.

Магистратура дает мне отсрочку, но не от реальной жизни, а от компромиссов.
Пока я студент, я защищен от необходимости продавать свой код ради еды. Это
позволяет направить всю энергию на создание того самого прототипа, который
станет моим пропуском в индустрию. Это не просто строчка в резюме. Это
доказательство того, что я умею доводить идею до состояния, пригодного для
демонстрации. Многие говорят, что игры — это командный вид спорта, но спортсмен
сначала тренируется в одиночестве, оттачивая мышечную память. Моя «мышечная
память» сейчас запоминает, как связывать анимации с механиками и как
оптимизировать сцены, не убивая FPS.

\section{Маленькая студия и большие надежды}

Когда диссертация будет защищена, и я официально выйду из академического
инкубатора, моя цель будет на первый взгляд скромной. Я не планирую штурмовать
штаб-квартиры гигантов индустрии с их многотысячными штатами и бесконечной
бюрократией. Мой путь лежит в маленькую студию, и я убежден, что география
здесь не имеет значения. Английский язык для меня — это не барьер, а дверь.
Рынок давно стал глобальным, и маленькая сплоченная команда где-нибудь в
Восточной Европе или Северной Америке для меня гораздо привлекательнее, чем
где либо.

Истоки классики, которую мы знаем, например, Fallout, создавались не в
конвейерных цехах, а в тесных офисах, где программист мог слышать дыхание
художника, а геймдизайнер — чертыхания композитора. Я хочу оказаться именно в
такой среде. В маленьких командах нет места иллюзиям. Здесь невозможно
спрятаться за спинами коллег. Если ты написал плохой код, игра рухнет на
ближайшем билде. Это жестоко, но именно так закаляется сталь профессионализма.

Моя цель на ближайшие два-три года после выпуска — стать «универсальным» в
рамках одного-двух завершенных проектов. Я хочу не просто написать модуль и
забыть о нем, а пройти весь цикл: от салфетки с наброском механики до момента,
когда игрок нажимает кнопку покупки в Steam. Мне нужно накопить опыт «закрытия»
проекта. Это навык, который трудно получить в одиночку. Успех в этот период я
измеряю не в деньгах, а в факте релиза. Проект, который дошел до пользователя и
получил отклик, успешен уже потому, что он состоялся. Это победа над энтропией.

\section{ Орден единомышленников }

К середине пятилетки, оглянувшись назад, я надеюсь увидеть рядом не просто
коллег, а людей, с которыми мы пережили краш-фиксы в три часа ночи и споры о
балансе. Именно из этих людей, проверенных огнем дедлайнов, я планирую собрать
ядро своей будущей компании. Небольшая команда близких друзей и
единомышленников.

В этот момент моя роль кардинально меняется. Я перестаю быть просто
программистом, наемным специалистом, которого нанимают решать типовые задачи. Я
становлюсь «одним из тех разработчиков». Тех, чьи игры обсуждают на форумах
любителей жанра. Моя специализация — иммерсивные симуляторы (immersive sims).
Это жанр, который требует от разработчика парадоксального мышления. Мы не пишем
сценарий, по которому игрок обязан пройти. Мы создаем физику, химию, логику
мира, а затем говорим игроку: «Ломай». Игры, где дверь можно не только взломать
отмычкой, но и вынести взрывчаткой, или пролезть через вентиляцию, или
прикинуться курьером — это высший пилотаж системного геймдизайна.

Моя будущая компания будет специализироваться именно на этом: на создании
миров, которые живут по своим законам, а не водят игрока за руку. Это нишевая
аудитория, но самая преданная и требовательная. Делать игры для них — значит
постоянно держать планку. Мы будем маленькой студией, но с большими амбициями.
В этой точке моя карьера становится неразрывно связанной с судьбой проекта.
Если наша игра «выстрелит», если имя студии станет синонимом качества в нише,
мы продолжим развивать это знамя.

\section{ Вовлеченность как компас }

Однако я прекрасно отдаю себе отчет в том, что индустрия непредсказуема. Проект
может оказаться гениальным, но не найденным аудиторией, или успешным финансово,
но выжигающим душу. Именно поэтому я не строю жестких планов на период «после
успеха или неуспеха». Вместо этого я закладываю в свою траекторию развилку,
основанную на философии Тима Кейна. Когда Кейн и его коллеги покинули
Interplay, они не просто искали работу — они искали смысл.

Если наша студия встанет на ноги, и мы почувствуем, что способны расти, не
теряя духа товарищества, мы продолжим путь независимых разработчиков. Но если
по какой-то причине формат собственной компании исчерпает себя, у меня есть
запасной, но не менее амбициозный путь — аутсорс-разработка высокого уровня.
Речь о продаже своего большого опыта и экспертизы тем проектам, которые мне
искренне нравятся. Это модель «свободного рыцаря» в мире геймдева. Такой подход
дает высочайший уровень контроля над своей жизнью: я выбираю проект не по
зарплатной ведомости, а по тому, загораются ли у меня глаза при чтении
дизайн-документа.

В основе всего этого лежит мое глубочайшее убеждение: вовлеченность — это один
из невосполнимых ресурсов, как и время. Деньги можно занять, репутацию —
восстановить. Но если у тебя нет внутреннего жжения при взгляде на рабочую
задачу, ты не создашь ничего стоящего. В иммерсив-симах нельзя обмануть игрока
фальшивым блеском, потому что суть этого жанра — интерактивность. Игрок
чувствует, если разработчику было скучно. Поэтому вовлеченность — это не хобби,
это двигатель моего прогресса, мой карьерный компас. Я не могу сказать точно,
какой именно код я буду писать через пять лет, но я знаю точно: этот код будет
частью игры, в которую я бы сам с удовольствем играл.

\end{document}
